\documentclass[10pt,a4paper]{amsart}
\usepackage[margin=19mm]{geometry}
\usepackage{amsmath,amssymb,mathtools}
\usepackage[scr=rsfs,cal=euler]{mathalfa}
\usepackage{xfrac}
\usepackage[unicode,hidelinks,bookmarks=false]{hyperref}
\newcommand{\ie}{i.e.\ }
\newcommand{\cf}{cf.\ }
\newcommand{\resp}{respectively\ }
\newcommand{\Subsection}{\subsection}
\providecommand{\nobreakdash}{\nobreak\hbox{-}}
\setlength{\emergencystretch}{2em}
\multlinegap0pt
\usepackage{amssymb}
\usepackage{mathtools}
\usepackage{enumerate}
\usepackage{dsfont}
\usepackage[normalem]{ulem}
\newtheorem{proposition}{Proposition}
\newcommand{\abs}[1]{\left|{#1}\right|}
\newcommand{\bb}{\mathbb}
\newcommand{\ds}{\mathrm{d}s}
\newcommand{\dt}{\mathrm{d}t}
\newcommand{\inpro}[2]{\left\langle#1,#2\right\rangle}
\newcommand{\norm}[1]{\left\|{#1}\right\|}
\title{Strong pathwise solutions for a class of stochastic thermo-magneto-hydrodynamic-type systems with multiplicative noise}
\author{Agus L. Soenjaya, Thanh Tran}
\date{Source: arXiv:2509.14490v3 (CC BY 4.0)}
\begin{document}
\maketitle

\noindent\emph{Source and licence.} Strong pathwise solutions for a class of stochastic thermo-magneto-hydrodynamic-type systems with multiplicative noise. Version arXiv:2509.14490v3 (CC BY 4.0), distributed by arXiv under the Creative Commons Attribution 4.0 International licence, http://creativecommons.org/licenses/by/4.0/, as recorded in the arXiv licence metadata for this exact version and shown on the abstract page https://arxiv.org/abs/2509.14490v3. Full licence notice: \texttt{licenses/arxiv-2509.14490-CC-BY-4.0.txt}.\par\smallskip

\noindent\textit{Editorial scope.} Counted unit: the article's proposition tau (source0.tex lines 1510--1520). The article's complete original proof of it is delivered with it (source0.tex lines 1523--1754, one \texttt{proof} environment of 232 lines). No mathematics is written, repaired or re-derived: the statement and the proof are the article's own lines, in the article's own notation, and no retained line is substituted or re-flowed.  Retained uncounted as the source-local context the proof needs: lines 254--262; lines 306--310; lines 1187--1189.  Nothing was omitted from any retained span, so the package carries no omission record.  No retained line contains a citation, so the wrapper carries no bibliography and no bibliography entry.  No retained line contains a figure environment, an image-inclusion command or drawing code, so no image is delivered and the package declares no asset dependency.  Editorial formatting only, in addition: an AMS \texttt{amsart} preamble that restates the article's own declarations and macros used by the retained lines, namely the environments \texttt{proposition} on separate counters, together with the standard packages those lines need.  The retained environments print in this document as Proposition 1 in order of appearance, whereas the article prints the same items with the numbering of its own full document.  The complete native source of the article as distributed in the arXiv e-print for this exact version is bundled unchanged as \texttt{sources/185248/source0.tex}.
	\begin{equation}\label{equ:est B with A}
		\abs{\mathcal{B}(v_1,v_2)}
		\leq
		\begin{cases}
			C \abs{v_1}^{\frac12} \norm{v_1}^{\frac12} \norm{v_2}^{\frac12} \abs{\mathcal{A}v_2}^{\frac12}, &\text{if $d=2$},
			\\[1ex]
			C \norm{v_1} \norm{v_2}^{\frac12} \abs{\mathcal{A} v_2}^{\frac12}, &\text{if $d\le 3$}.
		\end{cases}
	\end{equation}

    \begin{align}\label{equ:est R v in v}
		\abs{\mathcal{R}(v)}
		&\leq
		C \left(1+ \norm{v}^{r-1} \right) \norm{v}.
    \end{align}

\begin{align}\label{equ:Tn MT}
	\mathcal{T}_n^{M,T}= \left\{ \tau\leq T: \sup_{t\in [0,\tau]} \norm{\Phi^n(t)}^2 + \int_0^\tau \abs{\mathcal{A}\Phi^n(s)}^2 \ds  \leq 4M \right\}.
\end{align}

\begin{proposition}
Let $\mathcal{T}_n^{M,T}$ be the collection of stopping times defined in \eqref{equ:Tn MT}.
For any $n\in\bb{N}, \delta>0$, and stopping time $\tau$, define
\begin{align}\label{equ:event An tau}
	A_n(\tau,\delta):= \left\{ \sup_{t\in [0,\tau\wedge \delta]} \norm{\Phi^n(t)}^2 + \int_0^{\tau\wedge \delta} \abs{\mathcal{A}\Phi^n(s)}^2 \ds > 3M \right\}.
\end{align}
Then for any $T>0$,
\begin{align}\label{equ:lim sup prob}
	\lim_{\delta \to 0^+} \sup_{n\in\bb{N}}\, \sup_{\tau \in\mathcal{T}_n^{M,T}} \bb{P}\left[ A_n(\tau,\delta) \right] =0.
\end{align}
\end{proposition}

\begin{proof}
	Let $\tau\in\mathcal T_n^{M,T}$ and set
	\[
	\theta:=\tau\wedge\delta,
	\qquad
	I_n(t):=
	\int_0^t
	\abs{\mathcal A\Phi^n(s)}^2\,\mathrm ds .
	\]
	Since $\tau\in\mathcal T_n^{M,T}$, we have, for every $t\leq \theta$,
	\[
	\norm{\Phi^n(t)}^2\leq 4M \quad \text{and}
	\quad
	I_n(t)\leq 4M .
	\]
	Applying It\^o's formula to $\norm{\Phi^n}^2$ gives
	\begin{align}\label{equ:ito d Phi n small time}
		\mathrm{d} \norm{\Phi^{n}}^2
		+
		2 \abs{\mathcal{A} \Phi^{n}}^2 \dt
		&=
		\mathcal{D}_n(t)\,\dt
		+
		\mathrm d\mathcal M_n(t),
	\end{align}
	where
	\begin{align*}
	\mathcal{D}_n(t)
		&:=
		-2\inpro{\mathcal{B}(\Phi^n)}{\mathcal{A}\Phi^{n}}
		+
		2\inpro{\mathcal{R}(\Phi^n)}{\mathcal{A} \Phi^{n}}
		+
		\sum_{k=1}^\infty \norm{P_n g_k(\Phi^n)}^2,
		\\
		\mathcal M_n(t)
		&:=
		2\sum_{k=1}^\infty
		\int_0^t
		\inpro{g_k(\Phi^n(s))}{\mathcal A\Phi^n(s)}
		\,\mathrm d\beta_k(s).
	\end{align*}
	We estimate the drift term on the stopped interval $[0,\theta]$. By
	\eqref{equ:est B with A}, Young's inequality, and the bound
	$\norm{\Phi^n(t)}^2\leq 4M$, we have
	\[
	\abs{\inpro{\mathcal B(\Phi^n)}{\mathcal A\Phi^n}}
	\leq
	C\norm{\Phi^n}^{\frac32}
	\abs{\mathcal A\Phi^n}^{\frac32}
	\leq
	\frac14\abs{\mathcal A\Phi^n}^2+C_M .
	\]
	Similarly, by \eqref{equ:est R v in v} and Young's inequality,
	\[
	\abs{\inpro{\mathcal R(\Phi^n)}{\mathcal A\Phi^n}}
	\leq
	\frac14\abs{\mathcal A\Phi^n}^2+C_M .
	\]
	Moreover, by condition {\rm (G)},
	\[
	\sum_{k=1}^\infty \norm{P_n g_k(\Phi^n)}^2
	\leq
	C\left(1+\norm{\Phi^n}^2\right)
	\leq C_M .
	\]
	Consequently, for $t\leq\theta$, we obtain
	\[
	\mathcal{D}_n(t)\leq \abs{\mathcal A\Phi^n(t)}^2+C_M,
	\]
	where $C_M$ is independent of $n$, $\tau$, and $\delta$.
	
	Integrating \eqref{equ:ito d Phi n small time} from $0$ to $t\leq\theta$,
	we obtain
	\[
	\begin{aligned}
		\norm{\Phi^n(t)}^2+2I_n(t)
		&=
		\norm{\Phi_0^n}^2
		+
		\int_0^t \mathcal{D}_n(s)\,\mathrm ds
		+
		\mathcal M_n(t)
		\\
		&\leq
		\norm{\Phi_0^n}^2
		+
		I_n(t)
		+
		C_Mt
		+
		\mathcal M_n(t).
	\end{aligned}
	\]
	Hence, noting that $\norm{\Phi_0^n}^2\leq M$ a.s., we have
	\[
	\norm{\Phi^n(t)}^2+I_n(t)
	\leq
	M+C_M\delta+\mathcal M_n(t),
	\qquad \forall t\leq\theta .
	\]
	Therefore
	\begin{align}\label{equ:running-energy-bound}
		\sup_{t\in[0,\theta]}
		\left(
		\norm{\Phi^n(t)}^2+I_n(t)
		\right)
		\leq
		M+C_M\delta
		+
		\sup_{t\in[0,\theta]}\abs{\mathcal M_n(t)} .
	\end{align}
	
	Next, observe that
	\[
	\sup_{t\in[0,\theta]}\norm{\Phi^n(t)}^2
	+
	I_n(\theta)
	\leq
	2\sup_{t\in[0,\theta]}
	\left(
	\norm{\Phi^n(t)}^2+I_n(t)
	\right).
	\]
%	Indeed,
%	\[
%	\sup_{t\in[0,\theta]}\norm{\Phi^n(t)}^2
%	\leq
%	\sup_{t\in[0,\theta]}
%	\left(
%	\norm{\Phi^n(t)}^2+I_n(t)
%	\right),
%	\]
%	and also
%	\[
%	I_n(\theta)
%	\leq
%	\norm{\Phi^n(\theta)}^2+I_n(\theta)
%	\leq
%	\sup_{t\in[0,\theta]}
%	\left(
%	\norm{\Phi^n(t)}^2+I_n(t)
%	\right).
%	\]
	Thus, by \eqref{equ:running-energy-bound}, we infer that
	\[
	A_n(\tau,\delta)
	\subset
	\left\{
	M+C_M\delta
	+
	\sup_{t\in[0,\theta]}\abs{\mathcal M_n(t)}
	>
	\frac{3M}{2}
	\right\}.
	\]
	Choose $\delta_0>0$ sufficiently small so that $C_M\delta_0\leq \frac{M}{4}$.
	Then, for $\delta \in (0,\delta_0]$, we have
	\begin{align}\label{equ:An Mn delta}
	A_n(\tau,\delta)
	\subset
	\left\{
	\sup_{t\in[0,\theta]}\abs{\mathcal M_n(t)}
	>
	\frac{M}{4}
	\right\}.
	\end{align}
	
	It remains to estimate the martingale term $\mathcal{M}_n$. Its quadratic variation satisfies
	\begin{align*}
		\langle \mathcal M_n\rangle_\theta
		&=
		4\int_0^\theta
		\sum_{k=1}^\infty
		\abs{
			\inpro{g_k(\Phi^n(s))}{\mathcal A\Phi^n(s)}
		}^2
		\ds,
	\end{align*}
	where we recall that $\theta=\tau\wedge\delta$.
	By condition {\rm (G)} and the stopped bound $\norm{\Phi^n(s)}^2\leq 4M$, we obtain for all $s\leq\theta$,
	\[
	\sum_{k=1}^\infty
	\abs{
		\inpro{g_k(\Phi^n(s))}{\mathcal A\Phi^n(s)}
	}^2
	\leq
	C\norm{\Phi^n(s)}^2
	\sum_{k=1}^\infty \norm{g_k(\Phi^n(s))}^2
	\leq C_M.
	\]
	This implies $\langle \mathcal M_n\rangle_\theta
	\leq C_M\delta$.
	Thus, by the Burkholder--Davis--Gundy inequality and Chebyshev's inequality,
	\[
	\begin{aligned}
		\bb P\left[
		\sup_{t\in[0,\theta]}
		\abs{\mathcal M_n(t)}
		>
		\frac{M}{4}
		\right]
		\leq
		\frac{16}{M^2}
		\bb E\left[
		\sup_{t\in[0,\theta]}
		\abs{\mathcal M_n(t)}^2
		\right]
		\leq
		\frac{C}{M^2}
		\bb E\left[
		\langle\mathcal M_n\rangle_\theta
		\right]
		\leq
		C_M\delta .
	\end{aligned}
	\]
	Hence, for all $\delta\in (0,\delta_0]$, we obtain from \eqref{equ:An Mn delta},
	\[
	\bb P\left[A_n(\tau,\delta)\right]
	\leq C_M\delta,
	\]
	where $C_M$ is independent of $n$, $\tau$, and $\delta$. Taking the
	supremum over $n\in\bb N$ and $\tau\in\mathcal T_n^{M,T}$ yields
	\[
	\sup_{n\in\bb N}
	\sup_{\tau\in\mathcal T_n^{M,T}}
	\bb P\left[A_n(\tau,\delta)\right]
	\leq C_M\delta .
	\]
	Letting $\delta\to 0^+$ proves \eqref{equ:lim sup prob}.
\end{proof}

\end{document}
